% !TEX root = ../aa_SoM_Chapters.tex

%% HARRIS: Chapter 9
%\xdef\Isectiontitle{Statistical Mechanics} 
%%
%\xdef\IaSubsectiontitle{A Simple Thermodynamic System} 
%\xdef\IbSubsectiontitle{Entropy and Temperature}
%\xdef\IcSubsectiontitle{Einstein's solid}
%\xdef\IdSubsectiontitle{Boltzmann \& Maxwell–Boltzmann distributions}
%\xdef\IeSubsectiontitle{Classical Averages}
%
%\xdef\IfSubsectiontitle{Quantum Distributions} 
%\xdef\IgSubsectiontitle{The Quantum Gas} 
%\xdef\IhSubsectiontitle{Photon Gas}
%\xdef\IiSubsectiontitle{The Laser}
%\xdef\IjSubsectiontitle{Specific Heats} 
%\xdef\IkSubsectiontitle{Debye Model} 

% ö = \%c3\%b6
% ä = \%C3\%A4

\xdef\sectiontitle{\IIIsectiontitle} %aiheuttama

%%%%%%%%%%%%%%%
\begin{frame}[label=3_7_a]%[label=3_7_TOC1]
\frametitle{\texorpdfstring{\culine[\sectiontitleulinecolor]{\sectiontitle}}{\sectiontitle} }

\vspace{1cm}

\begin{itemize}[leftmargin=0.5cm,itemsep=3mm]
    \item[3.1] \hyperlink{3_1_a}{\IIIaSubsectiontitle}
    \item[3.2] \hyperlink{3_2_a}{\IIIbSubsectiontitle}	
    \item[3.3] \hyperlink{3_3_a}{\IIIcSubsectiontitle}	
    \item[3.4] \hyperlink{3_4_a}{\IIIdSubsectiontitle}
    \item[3.5] \hyperlink{3_5_a}{\IIIeSubsectiontitle}
    \item[3.6] \hyperlink{3_6_a}{\IIIfSubsectiontitle}
    \item[3.7] \hyperlink{3_7_a}{\CDAlert[blue]{\IIIgSubsectiontitle}}
\end{itemize}

%\endofslide 
\end{frame}
%%%%%%%%%%%%%%%

\xdef\sectiontitle{\IIIgSubsectiontitle} 


%%%%%%%%%%%%%%%
\begin{frame}
\frametitle{\texorpdfstring{\culine[\sectiontitleulinecolor]{\sectiontitle}}{\sectiontitle} }
\framesubtitle{NMR}

\begin{itemize}
	\item[] \CDAlert[fuchsia]{Nuclear Magnetic Resonance (NMR)}:
	
	\item Just as for \CDAlert[red]{Zeeman effect} for electrons in atoms\appear{, the \Alert[red]{energy levels for nuclei also split in external magnetic field}}

\item The \CDAlert[blue]{magnetic moment} for a nucleus is 
\[
\vec{\mu}  \equal \appear{\gamma} \appear{\vec{I}} \appear{\,, \qquad \mu_{N}}  \equal \frac{e \hbar}{2 m_{p}} \appear{\,, \qquad \text{where}} \qquad \appear{|\, \vec{I} \,|}  \equal  \appear{I(I+1)^{1 / 2} \hbar}
\]

\item The \CDAlert[fuchsia]{potential energy} for the magnetic moment in an external magnetic field is 
\[
U  \equal \appear{-\vec{\mu}} \appear{\cdot \vec{B}}
\]

\item The \CDAlert[fuchsia]{energy difference} between up and down states then becomes
\[
\Delta E  \equal \appear{2} \appear{( \mu_{z} )_{\text {proton}}} \appear{B}
\]
\end{itemize}

%\xdef\loc{south east}
%\begin{tikzpicture}[remember picture,overlay] % anchor=south
%    \node (A) [yshift=-0.25*\figYshift,xshift=\figXshift, anchor=\loc] at (current page.\loc) {\includegraphics[width=7cm]{Figures_booklet/SOM_12_Nuclear_physics_figures/SOM_Nuclear_physics_Figure_1.png}}; % 8cm max height
%\end{tikzpicture}

\endofslide 
%\endofsection
\end{frame}
%%%%%%%%%%%%%%%

%%%%%%%%%%%%%%%
\begin{frame}
\frametitle{\texorpdfstring{\culine[\sectiontitleulinecolor]{\sectiontitle}}{\sectiontitle} }
\framesubtitle{NMR}

\seminormal 
\begin{itemize}
	\item When hydrogen atoms \appear{(e.g. protons in water)} \appear{are \Alert[blue]{irradiated with electromagnetic radiation}}\appear{, i.e. photons of $h f=\Delta E$}\appear{, then we will have a \CDAlert[blue]{resonance absorption}, and \CDAlert[blue]{subsequent emission}}

\item In a magnetic field of $1 \mathrm{\;T}$\appear{, the energy difference $\Delta E=1.76 \times 10^{-7} \mathrm{\;eV}$} \appear{corresponds to resonance frequency $f=42.5$ $\mathrm{\;MHz}$}  \appear{(a radio frequency)}

\end{itemize}

% 6.5 cm = midway, 10 cm = 3/4 
\vspace{2.5mm}
\begin{minipage}{7.5cm}
\begin{itemize}
	\item NMR is applied to \Alert[fuchsia]{probe internal \CDAlert[fuchsia]{magnetic structure}} of materials\appear{, and especially in \Alert[fuchsia]{magnetic resonance imaging (\CDAlert[fuchsia]{MRI}) of living tissue}}\appear{, with negligible heating effect}
\appear{\xdef\loc{south east}
\begin{tikzpicture}[remember picture,overlay]% anchor=south
    \node (A) [yshift=-0.15*\figYshift,xshift=\figXshift, anchor=\loc] at (current page.\loc) {\includegraphics[width=5.8cm]{Figures_booklet/SOM_12_Nuclear_physics_figures/SOM_Nuclear_physics_Figure_17.png}};% 8cm max height
\end{tikzpicture}}
\item NMR is best known due to MRI\appear{, but it is widely applied also in other scientific areas}\appear{, such as \Alert[fuchsia]{observing the structures of} \CDAlert[fuchsia]{complex molecules in chemistry}}
\end{itemize}
\end{minipage}
%\vspace{2.5mm}

\endofslide 
%\endofsection
\end{frame}
%%%%%%%%%%%%%%%

%%%%%%%%%%%%%%%
\begin{frame}
\frametitle{\texorpdfstring{\culine[\sectiontitleulinecolor]{\sectiontitle}}{\sectiontitle} }
\framesubtitle{NAA}

% Neutron Activation Analysis (NAA)
\begin{itemize}[itemsep=2.2mm]
	\item[] \CDAlert[fuchsia]{Neutron Activation Analysis (NAA)}:
	\item A sample containing ${ }_{Z}^{A} M$ is exposed to a beam of slow neutrons

\item Nuclide ${ }^{A+1}_{~\;Z}M$ \appear{can be generated via reaction} $\appear{{ }_{Z}^{A} M} \appear{(}\appear{n,} \appear{\gamma)}\appear{^{A+1} M}$\appear{, but it \CDAlert[blue]{is radioactive}}\appear{, thus it will emit $\beta$ particles or $\gamma$ rays}

\item \CDAlert[blue]{Monitoring the emitted radiation}\appear{, ${ }^{A+1}_{~\;Z} M$ can be identified}\appear{, and furthermore concentration of ${ }_{Z}^{A} M $ in the sample can be determined}

\item This neatly enables \CDAlert[blue]{chemical analysis}\appear{, even different isotopes are distinguished from each other}\appear{, without breaking the sample}

\item The \CDAlert[fuchsia]{neutron activation analysis (NAA)} is suitable for analysing very small amounts of material\appear{, and often it is used 'when other physicochemical methods are not sensitive enough'}

\item Drawback is the difficult \CDAlert[red]{disposal} of radioactive samples

\end{itemize}

%\xdef\loc{south east}
%\begin{tikzpicture}[remember picture,overlay] % anchor=south
%    \node (A) [yshift=-0.25*\figYshift,xshift=\figXshift, anchor=\loc] at (current page.\loc) {\includegraphics[width=7cm]{Figures_booklet/SOM_12_Nuclear_physics_figures/SOM_Nuclear_physics_Figure_1.png}}; % 8cm max height
%\end{tikzpicture}

\endofslide 
%\endofsection
\end{frame}
%%%%%%%%%%%%%%%

%%%%%%%%%%%%%%%
\begin{frame}
\frametitle{\texorpdfstring{\culine[\sectiontitleulinecolor]{\sectiontitle}}{\sectiontitle} }
\framesubtitle{NAA}

\seminormal
\begin{itemize}
	\item After exposure to thermal neutrons the activity of a nuclide follows:
\[
R(t)  \equal \appear{\lambda} \appear{N(t)}  \equal \appear{R_{0}} \appear{(1-\rme^{-\lambda t} )}
\]
\appear{where $\lambda=$ decay constant} \appear{characteristic for ${ }^{A+1}_{~\;Z}M$}\appear{, and $R_{0}$ is the constant production rate of that isotope} \appear{($g$ in earlier Lecture)}

\item In the equation above\appear{, $R(t)$ is measured and $R_{0}$ is in principle computed from known values of intensity $I$} \appear{(specific for the set-up)} \appear{and cross section} \appear{(specific for the reaction leading to nuclei $ { }^{A+1}_{~\;Z}M$).} \appear{$I$ is the neutron flux (given e.g. in neutrons/s$\mathrm{m}^{2}$)}

\item In general the \CDAlert[blue]{rate for collisions/production} of is ${ }^{A+1}_{~\;Z}M$:
\[
R_{0}  \equal \appear{N_{0} \sigma I} \qquad \appear{, \qquad [\sigma I]}  \equal \appear{m^{2}} \frac{\text {neutrons}}{s \cdot m^{2}}  \equal \frac{\text {neutrons}}{s}
\]
\appear{where $N_{0}$ is the number of ${ }_{Z}^{A} M$ nuclei in the sample}
\end{itemize}

%\xdef\loc{south east}
%\begin{tikzpicture}[remember picture,overlay] % anchor=south
%    \node (A) [yshift=-0.25*\figYshift,xshift=\figXshift, anchor=\loc] at (current page.\loc) {\includegraphics[width=7cm]{Figures_booklet/SOM_12_Nuclear_physics_figures/SOM_Nuclear_physics_Figure_1.png}}; % 8cm max height
%\end{tikzpicture}

\endofslide 
%\endofsection
\end{frame}
%%%%%%%%%%%%%%%

%%%%%%%%%%%%%%%
\begin{frame}
\frametitle{\texorpdfstring{\culine[\sectiontitleulinecolor]{\sectiontitle}}{\sectiontitle} }
\framesubtitle{NAA}

\begin{itemize}
	\item Thereby, we obtain $\appear{R(t)}  \equal \appear{\lambda N(t)}  \equal \appear{N_{0} \sigma I (1-\rme^{-\lambda t} )}$ \appear{and if the half-life of the nuclide is short}\appear{, we obtain a saturation $R(t)=R(t=\infty)=N_{0} \sigma I$}

\item Table shows \CDAlert[red]{saturation activities $R(\infty)$} \appear{(per mg of sample)} \appear{for a few isotopes}
\appear{\xdef\loc{south east}%
\begin{tikzpicture}[remember picture,overlay]% anchor=south
    \node (A) [yshift=-0.25*\figYshift,xshift=\figXshift, anchor=\loc] at (current page.\loc) {\small \def\arraystretch{1.4}$
\begin{array}{c|c|c}
{ }_{Z}^{A} \mathrm{M} & { }^{A+1}_{~\;Z} \mathrm{M} & R(\infty) \text {\;decays} / \mathrm{min} \cdot \mathrm{mg} \\
\hline
{ }^{55} \mathrm{Mn} & { }^{56} \mathrm{Mn} & 8.8 \times 10^{6} \\
{ }^{63} \mathrm{Cu} & { }^{64} \mathrm{Cu} & 1.7 \times 10^{6} \\
{ }^{127} \mathrm{I} & { }^{128} \mathrm{I} & 1.6 \times 10^{6} \\
{ }^{197} \mathrm{Au} & { }^{198} \mathrm{Au} & 1.7 \times 10^{7}
\end{array}
$};% 8cm max height
\end{tikzpicture}}
\item Number of ${ }_{Z}^{A} M$ atoms in the sample is
\[
N_{0}  \equal \fraca{R(t=\infty)}{\sigma I} \appear{\,,}
\]
\appear{and mass of ${ }_{Z}^{A} M$ in sample is}
\end{itemize}

% 6.5 cm = midway, 10 cm = 3/4 
\vspace{2.5mm}
\begin{minipage}{7.8cm}
\begin{itemize}
	\item[] 
	\[
 \appear{m ({ }_{Z}^{A} M)}  \equal \frac{N_{0}}{N_{A}} \appear{M({ }_{Z}^{A} M)}  \equal \frac{R(t=\infty)}{N_{A} \sigma I} \appear{M({ }_{Z}^{A} M) \,,}
\]
\item[] where $M({ }_{Z}^{A} M)$ is the atomic weight of the element \appear{and $N_{A}$ is Avogadro's number}	
\end{itemize}
\end{minipage}
%\vspace{2.5mm}

%\xdef\loc{south east}
%\begin{tikzpicture}[remember picture,overlay] % anchor=south
%    \node (A) [yshift=-0.25*\figYshift,xshift=\figXshift, anchor=\loc] at (current page.\loc) {\includegraphics[width=7cm]{Figures_booklet/SOM_12_Nuclear_physics_figures/SOM_Nuclear_physics_Figure_1.png}}; % 8cm max height
%\end{tikzpicture}

\endofslide 
%\endofsection
\end{frame}
%%%%%%%%%%%%%%%

%%%%%%%%%%%%%%%
\begin{frame}
\frametitle{\texorpdfstring{\culine[\sectiontitleulinecolor]{\sectiontitle}}{\sectiontitle} }
\framesubtitle{Dating}

\begin{itemize}
	\item[] \CDAlert[fuchsia]{Radioactive Dating}:
\end{itemize}
	
%Radiaoctive Dating
\semismall
\begin{itemize}
	%\item[] \CDAlert[fuchsia]{Radiaoctive Dating}:
	\item Nature contains radioactive isotopes:
\begin{itemize}
	\item[(1)] Three heavy element $\alpha$-emitter chains \appear{(\CDAlert[blue]{thorium}}\appear{, \CDAlert[green]{uranium}}\appear{, \CDAlert[red]{actinium})}
\item[(2)] Isolated long-lived isotopes\appear{, such as ${ }^{40} \mathrm{K}$} \appear{with $T_{1 / 2}=1.25 \times 10^{9} \mathrm{\;y}$}
\item[(3)] Isolated radioactive isotopes are \CDAlert[fuchsia]{activated by cosmic rays}
\end{itemize}

\item For geological aging \appear{${ }^{40} \mathrm{K}$} \appear{and ${ }^{232} \mathrm{Th} ~(T_{1 / 2}=1.24 \times 10^{10} \mathrm{\;y})$} \appear{are used (for \CDAlert[blue]{"old" rocks}}\appear{, or "younger" rocks)}

\item \CDAlert[fuchsia]{Present abundance ratio} of two isotopes is measured\appear{, one of them being radioactive or a stable product of a radioactive decay\appear{}, (so that its abundance is}%time dependent)
\end{itemize}

% 6.5 cm = midway, 10 cm = 3/4 
\vspace{-0.3mm}
\begin{minipage}{5.7cm}
\begin{itemize}
	\item[] time dependent)\appear{, and compared against a \CDAlert[red]{previous abundance} \CDAlert[red]{ratio}}\appear{, that is known to have existed when the material was formed}
\end{itemize}
\end{minipage}
%\vspace{2.5mm}


\xdef\loc{south east}
\begin{tikzpicture}[remember picture,overlay] % anchor=south
    \node (A) [yshift=-0.25*\figYshift,xshift=\figXshift, anchor=\loc] at (current page.\loc) {\footnotesize \def\arraystretch{1.15}$
\begin{array}{c|c|c|c}
\text { Nuclide } & t_{1 / 2}(y) & \text { Abundance (\%) } & \text { Daughter } \\
\hline
{ }^{14} \mathrm{C} & 5730 & 1.35 \times 10^{-10} & { }^{14} \mathrm{N} \\
{ }^{40} \mathrm{K} & 1.25 \times 10^{9} & 0.0117 & { }^{40} \mathrm{A} \\
{ }^{87} \mathrm{Rb} & 4.88 \times 10^{10} & 27.83 & { }^{87} \mathrm{Sr} \\
{ }^{147} \mathrm{Sm} & 1.06 \times 10^{11} & 15.0 & { }^{143} \mathrm{Nd} \\
{ }^{176} \mathrm{Lu} & 3.59 \times 10^{10} & 2.59 & { }^{176} \mathrm{Hf} \\
{ }^{187} \mathrm{Re} & 4.30 \times 10^{10} & 62.60 & { }^{18} \mathrm{Os}
\end{array}
$}; % 8cm max height
\end{tikzpicture}

%\xdef\loc{south east}
%\begin{tikzpicture}[remember picture,overlay] % anchor=south
%    \node (A) [yshift=-0.25*\figYshift,xshift=\figXshift, anchor=\loc] at (current page.\loc) {\includegraphics[width=7cm]{Figures_booklet/SOM_12_Nuclear_physics_figures/SOM_Nuclear_physics_Figure_1.png}}; % 8cm max height
%\end{tikzpicture}

\endofslide 
%\endofsection
\end{frame}
%%%%%%%%%%%%%%%

%%%%%%%%%%%%%%%
\begin{frame}
\frametitle{\texorpdfstring{\culine[\sectiontitleulinecolor]{\sectiontitle}}{\sectiontitle} }
\framesubtitle{Dating}

\seminormal
\begin{itemize}
	\item In \CDAlert[red]{archeology}\appear{, ${ }^{14} \mathrm{C}$-dating is an important tool }
	\item ${ }^{14} \mathrm{C}$ is a $\beta^{-}$ emitter \appear{with half-life $T_{1 / 2}=5730$ years}

\item Due to the \CDAlert[blue]{carbon cycle in nature}\appear{, ratio of ${ }^{14} \mathrm{C} /{ }^{12} \mathrm{C}$ in living organisms equals the value in atmosphere} \appear{(presently the ratio ${ }^{14} \mathrm{C} /{ }^{12} \mathrm{C}=1.35 \times 10^{-12})$} 
\item When organism dies\appear{, compounds such as atmospheric $\mathrm{CO}_{2}$} \appear{or hydrocarbons} \appear{stop circulating through the organism} 
\implies{ Carbon stops being renewed\appear{, and the radioactive ${ }^{14} \mathrm{C}$ starts to decay}}

\item \CDAlert[red]{Activity} \appear{(decay rate)} \appear{of one gram of carbon in the material} \appear{allows to calculate the \CDAlert[red]{time of death} of the organism.} \appear{In reality ${ }^{14} \mathrm{C} /{ }^{12} \mathrm{C}$ fraction} \appear{has been varying through years}\appear{, which weakens the accuracy of dating process} \vspace{2mm}

\begin{itemize} \small
	\item ${\sim}$9000 years ago the ratio was $1.5$ times the ratio of the present time
	\item Nuclear experiments have increased the ratio
	\item \Alert[fuchsia]{Use of fossil fuel has decreased the ratio} \appear{(signs of the \CDAlert[fuchsia]{Anthropocene?})}
\end{itemize}

\end{itemize}

%\xdef\loc{south east}
%\begin{tikzpicture}[remember picture,overlay] % anchor=south
%    \node (A) [yshift=-0.25*\figYshift,xshift=\figXshift, anchor=\loc] at (current page.\loc) {\includegraphics[width=7cm]{Figures_booklet/SOM_12_Nuclear_physics_figures/SOM_Nuclear_physics_Figure_1.png}}; % 8cm max height
%\end{tikzpicture}

\endofslide 
%\endofsection
\end{frame}
%%%%%%%%%%%%%%%

%%%%%%%%%%%%%%%
\begin{frame}
\frametitle{\texorpdfstring{\culine[\sectiontitleulinecolor]{\sectiontitle}}{\sectiontitle} }
\framesubtitle{Dating}

\begin{itemize}
	\item Let's find out how \CDAlert[blue]{the age of the Earth} can be estimated!

\item In radioactive decay, where a parent $(P)$ nuclide \appear{decays into daughter $(D)$}\appear{, the numbers of parent and daughter nuclei are respectively given by:} 
\[
N_{P}(t)  \equal \appear{N_{0}} \appear{\rme^{-\lambda t}} \qquad \appear{, \qquad N_{D}(t)}  \equal \appear{N_{0}-N_{P}}
\]

\item Then, at each given moment of time $t$, we have 
\[
t \equal \frac{1}{\lambda} \appear{\ln} \LP \frac{N_{0}}{N_{P}} \;\RP  \equal \frac{T_{1 / 2}}{\ln 2} \appear{\ln} \LP \frac{N_{0}}{N_{P}} \,\RP  \equal \frac{T_{1 / 2}}{\ln 2} \appear{\ln} \LP \appear{1} \plus \frac{N_{D}}{N_{P}} \,\RP
\]

\item Some isotope ratios \appear{(}$\appear{{ }^{238} \mathrm{U} /{ }^{206} \mathrm{Pb}}\appear{,~{ }^{87} \mathrm{Rb} /{ }^{87} \mathrm{Sr}}\appear{,~{ }^{40} \mathrm{K} /{ }^{40} \mathrm{Ar}}$\appear{)} \appear{are used as 'rock-clocks'}\appear{, e.g. indicating that the \CDAlert[blue]{age of Earth's surface rock} \CDAlert[blue]{is around $4.5 \times 10^{9}$ years}}
\end{itemize}

%\xdef\loc{south east}
%\begin{tikzpicture}[remember picture,overlay] % anchor=south
%    \node (A) [yshift=-0.25*\figYshift,xshift=\figXshift, anchor=\loc] at (current page.\loc) {\includegraphics[width=7cm]{Figures_booklet/SOM_12_Nuclear_physics_figures/SOM_Nuclear_physics_Figure_1.png}}; % 8cm max height
%\end{tikzpicture}

\endofslide 
%\endofsection
\end{frame}
%%%%%%%%%%%%%%%

%%%%%%%%%%%%%%%
\begin{frame}
\frametitle{\texorpdfstring{\culine[\sectiontitleulinecolor]{\sectiontitle}}{\sectiontitle} }
\framesubtitle{Alarms}

% Radioactive nuclei in fire alarms
\begin{itemize}
	\item Am-241 is an isotope belonging to Neptunium series 
	
	\item Am-241 is used in conventional household fire alarms to ionize the air

\item \CDAlert[red]{In clean air}, the ions in air move easily to electrodes \appear{and create a small trickle current }
\implies{ this \CDAlert[red]{trickle current is continuously measured} }

\item \CDAlert[blue]{If smoke particles enter} the channel\appear{, they remove the ions (of the ionized gas molecules) from the air}\appear{, electrical conductivity of the air decreases} 
\implies{ the trickle current stops }
\end{itemize}

%\xdef\loc{south east}
%\begin{tikzpicture}[remember picture,overlay] % anchor=south
%    \node (A) [yshift=-0.25*\figYshift,xshift=\figXshift, anchor=\loc] at (current page.\loc) {\includegraphics[width=7cm]{Figures_booklet/SOM_12_Nuclear_physics_figures/SOM_Nuclear_physics_Figure_1.png}}; % 8cm max height
%\end{tikzpicture}

\endofslide 
%\endofsection
\end{frame}
%%%%%%%%%%%%%%%

%%%%%%%%%%%%%%%
\begin{frame}
\frametitle{\texorpdfstring{\culine[\sectiontitleulinecolor]{\sectiontitle}}{\sectiontitle} }

\begin{itemize}
	\item Particle-Induced \CDAlert[red]{X-ray} and \CDAlert[blue]{Gamma-ray} Emission (PIXE and PIGE)
\item \CDAlert[red]{PIXE}: 
\item[] \begin{itemize}
	 \item An elemental analysis of \CDAlert[red]{trace element concentrations} of samples \appear{(for elements $Z>20$)} 
	 \item Sample is bombarded with low energy ions \appear{(few MeV)}\appear{, \CDAlert[red]{X-rays} will be emitted} \appear{\CDAlert[red]{from K- and L-shell electrons} of the atoms}
	 
	 \item This offers a method to identify elements in the sample 
	 
	 \item Depending on the ions used\appear{, the method \Alert[fuchsia]{can be very surface-sensitive} }
	 
	 \item PIXE can be used e.g. to detect elemental composition in trees \appear{(which have the annual growth rings)} \appear{or e.g. verify genuinity of old paintings} \appear{(with distinct elemental composition)}
\end{itemize}

%\item \CDAlert[blue]{PIGE}: Bombarded with higher ion energies, nuclei start emitting $\gamma$-rays from the nucleus in a similar way. Then also lighter elements can be detected (PIGE = Particle-Induced Gamma-ray Emission). For example, the ${ }^{23} \mathrm{Na}\left(\mathrm{p}, \mathrm{p}^{\prime} \mathrm{v}\right)^{23} \mathrm{Na}$ reaction is frequently used for determining bulk sodium content of glasses 
%
%\item PIXE and PIGE are used widely in elemental analysis in material and environmental sciences.
\end{itemize}

%\xdef\loc{south east}
%\begin{tikzpicture}[remember picture,overlay] % anchor=south
%    \node (A) [yshift=-0.25*\figYshift,xshift=\figXshift, anchor=\loc] at (current page.\loc) {\includegraphics[width=7cm]{Figures_booklet/SOM_12_Nuclear_physics_figures/SOM_Nuclear_physics_Figure_1.png}}; % 8cm max height
%\end{tikzpicture}

\endofslide 
%\endofsection
\end{frame}
%%%%%%%%%%%%%%%

%%%%%%%%%%%%%%%
\begin{frame}
\frametitle{\texorpdfstring{\culine[\sectiontitleulinecolor]{\sectiontitle}}{\sectiontitle} }

\begin{itemize}
	\item Particle-Induced \CDAlert[red]{X-ray} and \CDAlert[blue]{Gamma-ray} Emission (PIXE and PIGE)

	\item \CDAlert[blue]{PIGE}: 
	\item[] \begin{itemize}
	 \item Bombarded with higher ion energies\appear{, nuclei start emitting \Alert[blue]{$\gamma$ rays from the nucleus}} \appear{in a similar way to PIXE}
	 \item Then also lighter elements can be detected 
	  
	 \item For example, the $\appear{{ }^{23} \mathrm{Na}}\appear{( \mathrm{p}}\appear{,} \appear{\mathrm{p}^{\prime}} \appear{+ \gamma )} \appear{{}^{23} \mathrm{Na}}$ \appear{reaction is frequently used for \CDAlert[blue]{determining bulk sodium content of glasses}} 
\end{itemize}

\item[]
\item PIXE and PIGE are both used widely in \CDAlert[fuchsia]{elemental analysis in} \appear{\CDAlert[fuchsia]{material and environmental sciences}}
\end{itemize}

%\xdef\loc{south east}
%\begin{tikzpicture}[remember picture,overlay] % anchor=south
%    \node (A) [yshift=-0.25*\figYshift,xshift=\figXshift, anchor=\loc] at (current page.\loc) {\includegraphics[width=7cm]{Figures_booklet/SOM_12_Nuclear_physics_figures/SOM_Nuclear_physics_Figure_1.png}}; % 8cm max height
%\end{tikzpicture}

%\endofslide 
\endofsection
\end{frame}
%%%%%%%%%%%%%%%

